% Cote 119: offprints, reports, courses, CV and letters (typed and printed part): archive cote 119, site page 229 (PDF page 230).
% typescript. Draft compared: none.
% AI-transcribed by Claude Opus 5.5 (claude-opus-5-5, adaptive thinking, effort high), Message Batches API, batch r119_full (msgbatch_01DwZNm5Y7YsKBWyAN4sFs8w), 2026-11-13; not yet checked by eye.
\scanpage{119/229}
\mspage{. 16 .}

g) \uline{Théorie des intersections et théorème de Riemann-Roch} :

La principale idée nouvelle, c'est qu'il y
a presque identité entre le groupe ``de Chow'' des
classes de cycles sur une variété $X$ , et un certain
groupe de ``classes de faisceaux cohérents'' (tout au
moins modulo torsion), à savoir le groupe $K(X)$
(mentionné dans le par. 3). Dans un contexte modeste
c'est exposé dans [12] et le travail de A. Borel et
J.P. Serre [$*$], dans un contexte plus ambitieux cela
donne l'imposant séminaire [16, SGA 7]. Dans le même
esprit, cf. [12 bis].

Par ailleurs, l'idée (que je semble avoir été
le premier à introduire avec ma formulation du théorème
de Riemann-Roch) de reformuler un théorème sur une
variété (dû en l'occurrence à F. Hirzebruch) en un
théorème plus général sur un morphisme de variétés, a
connu par la suite une grande fortune, non seulement
en géométrie algébrique, mais aussi en topologie
algébrique et topologie différentielle (à commencer par
la ``formule de Riemann-Roch différentiable'', développée
par M.F. Atiyah et F. Hirzebruch sous l'inspiration de
ma formulation ``relative'' du théorème de Riemann-Roch).

h) \uline{Schémas abéliens} :

En termes plus classiques, ce sont les familles
de variétés abéliennes, paramétrées par un schéma
quelconque. Les résultats les plus importants que j'y
ai établis sont le ``\uline{théorème de réduction semi-stable}''
et ses conséquences et variantes [16, SGA 7, Exp. IX],
le théorème d'\uline{existence de morphismes de schémas abéliens}
contenu dans [17] et ses variantes (généralisé par
P. Deligne [$*$] en un théorème sur la cohomologie de

\margin{229}
